\scanpage{5}
Noter que $\mathrm{cd}^{\text{ét}}_p(k)$ est égal à $0$ ou $1$, suivant
que \deleted{toute} extension \inserted{\figure{\begin{tikzpicture}[baseline=(n.base)]\node[draw,ellipse,inner sep=1pt] (n) {finie};\end{tikzpicture}}} étale de $k$ est d'ordre
premier à $p$, ou non.

Si $G$ est de t. m., alors $G$ est dual
d'un groupe fini étale $C$ lequel est
quotient d'un module galoisien, \unsure{libre} de
type fini sur $\mathbf{Z}$, donc $G$ est contenu
dans un tore, donc on a une suite
exacte
\diagram{\begin{tikzcd}
0\arrow[r]&G\arrow[r]&T\arrow[r]&T'\arrow[r]&0
\end{tikzcd}}
$T$ et $T'$ des tores. D'où le

\uline{Lemme 2} Pour que $\mathrm{cd}^{\text{t.m.}}_p(k)$ soit $\leq n$,
\margin{$\left[\rule{0pt}{6em}\right.$}
il f. et s. que les conditions suivantes
soient vérifiées: pour tt tore $T$ sur $k$,
$H^i(k,T)(p)$ est nul pour $i>n$, et $p$-divisible
si $i=n$.
\inserted{[Ou encore: pour tte extension finie \unsure{sép.} $k'$ de $k$ [\unsure{il suffit} d'ordre premier à $p$]
\deleted{$\operatorname{Br}(k)(p)$} $H^i(k',G_m)=0$ \inserted{si $i>n$, et $H^{\unsure{n}}(k',G_m)$ est $p$-divisible}]}

Comme on a
\diagram{\begin{tikzcd}
0\arrow[r]&{}_{p}T\arrow[r]&T\arrow[r,"p"]&T\arrow[r]&0
\end{tikzcd}}
d'où
\[
\deleted{\left\{ {}_{p}H^{\unsure{i|n}}(T)=\operatorname{Im}H^i({}_{p}T)\right.}
\]
\margin{4}

